\section*{第 \ref{ch:g7:literal} 章 --- 代数式}

\begin{solution}{exo:g7:literal:1}
$7x$; \quad $3y$; \quad $5ab$; \quad $x^2$; \quad $2(a + 4)$.
\end{solution}

\begin{solution}{exo:g7:literal:2}
$5ab = 5 \times a \times b$; \quad
$3(x+2) = 3 \times (x + 2)$; \quad
$x^2 + 4x = x \times x + 4 \times x$; \quad
$2a^3 = 2 \times a \times a \times a$.
\end{solution}

\begin{solution}{exo:g7:literal:3}
当 $x = 3$ 时：\ $5 \times 3 + 1 = 16$；\quad $3 \times 3 = 9$；\quad
$2 \times 9 - 3 = 15$; \quad $10 - 6 = 4$.
\end{solution}

\begin{solution}{exo:g7:literal:4}
$3a + 2b = 12 + 3 = 15$; \qquad
$ab + 5 = 4 \times 1.5 + 5 = 6 + 5 = 11$.
\end{solution}

\begin{solution}{exo:g7:literal:5}
$8x + 5x = 13x$; \quad
$9a - 4a + a = 6a$; \quad
$6y + 2 + 3y + 7 = 9y + 9$; \quad
$5t + 3 - 2t - 3 = 3t$.
\end{solution}

\begin{solution}{exo:g7:literal:6}
正方形：$P = 4c$。等腰三角形：$P = b + 2s$。面包的价格：
$p = 1.20\,n$ 欧元。
\end{solution}

\begin{solution}{exo:g7:literal:7}
$2n$：两倍；\ $n + 2$：比 $n$ 多 2；\ $n^2$：平方；
\ $\frac n2$：一半。
\end{solution}

\begin{solution}{exo:g7:literal:8}
$3(x+4) = 3x + 12$：\emph{成立}（对任意 $x$，依据分配律 ---
这是一个证明，无需代入检验）。

$2x + 5 = 7x$：取 $x = 2$ 检验：左边为 $9$，右边为 $14$ ---
\emph{一般不成立}（只在 $x = 1$ 时成立）。

$x + x = x^2$：取 $x = 3$ 检验：左边为 $6$，右边为 $9$ --- \emph{不成立}
（一般情况下 $x + x = 2x$）。
\end{solution}

\begin{solution}{exo:g7:literal:9}
\emph{1.} 长：$w + 3$。周长：
\[
P = 2\bigl(w + (w + 3)\bigr) = 2(2w + 3) = 4w + 6 .
\]

\emph{2.} 当 $w = 5$ 时：直接看，两边长分别为 $5$ 和 $8$，于是
$P = 2 \times 13 = 26$ cm；由公式得 $4 \times 5 + 6 = 26$ cm。
两者一致。
\end{solution}

\begin{solution}{exo:g7:literal:10}
按字母 $n$ 进行以下步骤：
\[
n \ \to\ n + 6 \ \to\ 3(n + 6) = 3n + 18 \ \to\ 3n + 18 - 18 = 3n
\ \to\ \frac{3n}{n} = 3 .
\]
所有人得到的结果都是 $3$。
\end{solution}

\begin{solution}{exo:g7:literal:11}
写成差的形式：$A = a^2 - b^2$。将 L 形切成两个
长方形：一条全宽的长条，高为 $a - b$，上面再放一个
较窄的长方形：
\[
A = a(a - b) + b(a - b)
\]
（下方长方形宽 $a$、高 $a - b$，上方长方形宽 $a - b$
、高 $b$，得到另一种写法 $a(a-b) + (a-b)b$ --- 结果相同）。
当 $a = 5$，$b = 2$ 时：$a^2 - b^2 = 25 - 4 = 21$，且
$5 \times 3 + 2 \times 3 = 15 + 6 = 21$。两个表达式结果一致 ---
事实上它们总是相等：这就是 $a^2 - b^2 = (a + b)(a - b)$ 这一恒等式，
只是以另一种形式出现在\cref{ch:g9:algebra}中。
\end{solution}

\begin{solution}{pb:g7:literal:1}
\textbf{1.} 图形 $1$ 至 $4$ 分别需要 $4$、$7$、$10$、$13$
根火柴。

\textbf{2.} 每新增一个正方形，它与前一图形共用一条边，
因此只需 $3$ 根新火柴（共 $4$ 条边）。由
图形 $4$（$13$ 根火柴）再加六个正方形，需要
$6 \times 3 = 18$ 根，图形 $10$ 需要 $13 + 18 = 31$ 根火柴。

\textbf{3.} 从左往右看：先是一根竖火柴，
然后 $n$ 个正方形中的每一个各贡献一条
上边、一条下边和一条右边——每个 $3$ 根。总数：
$3n + 1$；“$+1$”正是最开头那根竖火柴。验证：
$n = 1, 2, 3, 4$ 分别给出 $4, 7, 10, 13$。

\textbf{4.} $4 + 3(n - 1) = 4 + 3n - 3 = 3n + 1$：阿米尔的
公式展开后恰好得到同样的化简式。完全
一致。

\textbf{5.} 莉娜：$2n + (n + 1) = 3n + 1$ --- 结果又是一样的。
三种看法，同一个化简式。图形 $100$：
需要 $3 \times 100 + 1 = 301$ 根火柴。

\textbf{6.} 火柴数依次为 $3$、$5$、$7$ 根。每个新增三角形
借助一条已有的边，只增加 $2$ 根火柴，
从 $3$ 根开始：图形 $n$ 需要 $3 + 2(n - 1) = 2n + 1$ 根火柴。

\textbf{7.} 逆推 $2n + 1 = 49$：去掉第一根火柴后剩下
$48$，此后每个三角形占用 $2$ 根：
$48 \div 2 = 24$ 个三角形。

\textbf{8.} 排成一排的 $n$ 张桌子，每张桌子提供上、下两条边
的座位（共 $2n$ 个），两端的两张桌子各多出一条边
的座位（各 $+2$ 个）：共 $2n + 2$ 个座位。检验：$4$、$6$、$8$
分别对应 $n = 1, 2, 3$ 时的情形。

\textbf{9.} 排成一行时：$2n + 2 \geq 30$，即 $2n \geq 28$，所以
需要 $14$ 张桌子。将同样的 $14$ 张桌子分成两行，每行
$7$ 张：每行有 $2 \times 7 + 2 = 16$ 个座位，合计 $32$ 个 ---
比单行多两个，因为每增加一行就多出两个新的
\emph{两端}。（每次分行多出 $2$ 个座位，宴会承办人对此心知肚明。）

\textbf{10.} 大小为 $n$ 的 L 形，竖向有 $n$ 个点，横向有 $n$ 个点，
拐角处的点共用：共 $n + n - 1 = 2n - 1$ 个点。
图形 $1$、$2$、$3$ 分别有 $1$、$3$、$5$ 个点 --- \emph{奇
数}。将这些 L 形一个套一个嵌套，恰好
填满一个 $n \times n$ 的点阵正方形：这正是
\cref{pb:g8:equations:1}中的故事。

\textbf{11.} $3(n + 2) - 5 = 3n + 6 - 5 = 3n + 1$：该断言
化为已证公式，因此对每个 $n$ 都成立 ---
这是代数推出的结论，而非由例子得来。

\textbf{12.} 检验：$4 \geq 2$、$9 \geq 3$、$100 \geq 10$：全部
成立。但 $n = 0.5$ 时 $n^2 = 0.25 < 0.5$：该断言不成立。
三次检验并没有证明关于\emph{所有}数的任何结论（它们
只是未能发现问题）；唯一一个反例
\emph{证明}了该一般断言不成立
（见\cref{met:g7:literal:test}）。字母代表每一个数
--- 包括小数。

\textbf{13.} 区域数：$2$ 个点 $\to 2$；$3$ 个点 $\to 4$；
$4$ 个点 $\to 8$；$5$ 个点 $\to 16$。规律似乎暗示
“每增加一个点，区域数翻一倍”，预期有 $32$ 个区域，
对应 $6$ 个点。

\textbf{14.} 仔细计数得到 $31$ 个区域，而不是 $32$。
翻倍猜想已经破灭：它经受住了前四次检验，
却在第五次检验中失败。（这一反例有名到值得专门一个名字，
称为“圆内区域陷阱”；等间距取点之类的捷径并不可行 ---
若三条弦交于同一点，区域会合并，
从而破坏计数。）

\textbf{15.} 证明需要一个\emph{结构性}的理由，解释为何加入一个点
会使区域数翻倍 ---但这样的理由根本不存在。
翻倍只是小情形下的巧合。火柴公式之所以可靠，
是因为每个公式都直接从作图本身读出
（每个正方形明显占用 $3$ 根火柴，
每张桌子明显提供 $2$ 个座位）：靠的是结构，而不是检验。
已证明的预测：$3 \times 2\,026 + 1 = 6\,079$ 根火柴。
\end{solution}
